% !TEX root = ../thesis.tex

Legged robots equipped with robotic arms can perform complex loco-manipulation
tasks that neither locomotion-only nor fixed-base manipulation systems can
achieve alone.
This project adapts the open-source \textit{RoboDuet} framework---originally
developed at Shanghai Jiao Tong University for the Unitree Go1+Z1
platform---to the heavier and more capable \textit{Unitree B2+Z1} system,
which comprises a 60\,kg quadruped body paired with a six-degrees-of-freedom
robotic arm.

The core technical contribution is a complete migration of the decentralised
multi-agent reinforcement learning (MARL) pipeline: robot URDF integration,
degree-of-freedom (DOF) mapping, PD gain calibration, and the development of
a simulation demonstration harness using NVIDIA IsaacGym.
Two decoupled policies are trained via Proximal Policy Optimisation (PPO):
a \textit{locomotion policy} that stabilises the quadruped at a nominal
standing height ($h\approx0.45$\,m), and an \textit{arm policy} that tracks
spherical end-effector commands $(l, p, \psi)$---reach distance, pitch, and
yaw---using a history-conditioned actor-critic architecture with a learned
adaptation module for sim-to-real transfer.

Three demonstrations are produced in simulation:
(1)~fixed-base arm manipulation confirming forward end-effector reach;
(2)~standing quadruped with simultaneous arm reaching; and
(3)~a fixed-base grasp-and-lift sequence in which a small 6\,cm cube is
picked up and elevated by 30\,cm, demonstrating a complete reach-grasp-lift
pipeline using kinematic box teleportation.
The robot base height remains within $\pm$3\,cm of nominal in Demonstration~2,
and is kinematically fixed at 0.340\,m in Demonstrations~1 and~3, confirming
the stability of the concurrent dog-arm control in the standing case.

These results confirm the cross-morphology zero-shot transfer capability
claimed in the original RoboDuet paper and establish the B2Z1 platform as
a viable testbed for future real-world loco-manipulation experiments.
